% test_plan.tex — Test Plan per la funzionalità "Gestione Bug"

\section{Test Plan – Gestione Bug}

\subsection{Obiettivo}
Il test plan descrive le attività di verifica previste per la funzionalità
\textbf{Gestione Bug} del sistema BugBoard26.
L'obiettivo è validare che le operazioni di creazione, assegnazione e
modifica di un bug rispettino i requisiti funzionali e le regole di
autorizzazione definite nella specifica.

\subsection{Funzionalità sotto test}
Le funzionalità coperte dal piano sono quelle esposte dal servizio
\texttt{BugService}:

\begin{enumerate}
    \item \textbf{Creazione di un bug} – metodo \texttt{create(BugCreateRequest, UUID)}
    \item \textbf{Assegnazione di un bug} – metodo \texttt{assign(UUID, AssignRequest, UUID)}
    \item \textbf{Modifica di un bug} – metodo \texttt{patch(UUID, BugPatchRequest, UUID, boolean)}
\end{enumerate}

\subsection{Approccio di testing}

\begin{itemize}
    \item \textbf{Livello}: test di unità (\emph{unit test}).
    \item \textbf{Tecnica}: \emph{black-box} (classi di equivalenza,
          valori limite) combinata con ispezione \emph{white-box}
          (copertura dei rami decisionali interni).
    \item \textbf{Framework}: JUnit~5 + Mockito~5 (via
          \texttt{spring-boot-starter-test}).
    \item \textbf{Isolamento}: le dipendenze (repository JPA) sono sostituite
          da \emph{mock objects} iniettati con \texttt{@Mock} /
          \texttt{@InjectMocks}, garantendo che i test verifichino
          esclusivamente la logica di business senza accesso al database.
\end{itemize}

\subsection{Casi di test pianificati}

\subsubsection{create(BugCreateRequest, UUID)}

\begin{center}
\small
\begin{tabular}{|c|p{6cm}|p{5.5cm}|}
\hline
\textbf{ID} & \textbf{Scenario} & \textbf{Risultato atteso} \\
\hline
TC-C1 & Creazione con titolo, descrizione, tipo e priorità validi & Bug salvato; history \texttt{CREATE} registrata \\
\hline
TC-C2 & UserId inesistente nel sistema & \texttt{EntityNotFoundException} \\
\hline
TC-C3 & Request con lista di label (alcune esistenti, altre nuove) & Bug con label associate; label inesistenti create automaticamente \\
\hline
\end{tabular}
\end{center}

\subsubsection{assign(UUID, AssignRequest, UUID)}

\begin{center}
\small
\begin{tabular}{|c|p{6cm}|p{5.5cm}|}
\hline
\textbf{ID} & \textbf{Scenario} & \textbf{Risultato atteso} \\
\hline
TC-A1 & Admin assegna bug a utente esistente & Assignee aggiornato; history \texttt{ASSIGN} registrata \\
\hline
TC-A2 & Utente non-admin tenta l'assegnazione & \texttt{SecurityException} \\
\hline
TC-A3 & BugId inesistente & \texttt{EntityNotFoundException} \\
\hline
\end{tabular}
\end{center}

\subsubsection{patch(UUID, BugPatchRequest, UUID, boolean)}

\begin{center}
\small
\begin{tabular}{|c|p{6cm}|p{5.5cm}|}
\hline
\textbf{ID} & \textbf{Scenario} & \textbf{Risultato atteso} \\
\hline
TC-P1 & Admin aggiorna titolo e priorità & Campi aggiornati; history \texttt{UPDATE} registrata \\
\hline
TC-P2 & Non-admin tenta di cambiare stato di bug altrui & \texttt{SecurityException} \\
\hline
TC-P3 & Assegnatario modifica il proprio bug senza cambiare stato & Modifica applicata con successo \\
\hline
\end{tabular}
\end{center}

\subsection{Criteri di ingresso e uscita}

\begin{description}
    \item[Ingresso:] codice sorgente compilabile; dipendenze Maven risolte;
    ambiente Java~23 configurato.
    \item[Uscita:] tutti i 9 test passano (\texttt{BUILD SUCCESS}) con
    0~failures e 0~errors; copertura dei rami decisionali principali
    raggiunta.
\end{description}

\subsection{Ambiente di esecuzione}

\begin{itemize}
    \item \textbf{JDK}: Amazon Corretto 23.0.2
    \item \textbf{Build tool}: Apache Maven con plugin Surefire 3.1.2
    \item \textbf{CI locale}: \texttt{mvn test} dalla directory \texttt{backend/}
\end{itemize}
